\documentclass[fleqn,10pt]{latex/stylish_article} % Document font size and equations flushed left

\setcounter{tocdepth}{3}

% Pandoc environments
\usepackage{framed}
\usepackage{fancyvrb}
\providecommand{\tightlist}{%
  \setlength{\itemsep}{0pt}\setlength{\parskip}{0pt}}
\newcommand{\VerbBar}{|}
\newcommand{\VERB}{\Verb[commandchars=\\\{\}]}
\DefineVerbatimEnvironment{Highlighting}{Verbatim}{commandchars=\\\{\}, fontsize=\scriptsize} % R Code

% Colored code
\usepackage{color}
\definecolor{shadecolor}{RGB}{248,248,248}
\newenvironment{Shaded}{\begin{snugshade}}{\end{snugshade}}
\newcommand{\KeywordTok}[1]{\textcolor[rgb]{0.13,0.29,0.53}{\textbf{{#1}}}}
\newcommand{\DataTypeTok}[1]{\textcolor[rgb]{0.13,0.29,0.53}{{#1}}}
\newcommand{\DecValTok}[1]{\textcolor[rgb]{0.00,0.00,0.81}{{#1}}}
\newcommand{\BaseNTok}[1]{\textcolor[rgb]{0.00,0.00,0.81}{{#1}}}
\newcommand{\FloatTok}[1]{\textcolor[rgb]{0.00,0.00,0.81}{{#1}}}
\newcommand{\ConstantTok}[1]{\textcolor[rgb]{0.00,0.00,0.00}{{#1}}}
\newcommand{\CharTok}[1]{\textcolor[rgb]{0.31,0.60,0.02}{{#1}}}
\newcommand{\SpecialCharTok}[1]{\textcolor[rgb]{0.00,0.00,0.00}{{#1}}}
\newcommand{\StringTok}[1]{\textcolor[rgb]{0.31,0.60,0.02}{{#1}}}
\newcommand{\VerbatimStringTok}[1]{\textcolor[rgb]{0.31,0.60,0.02}{{#1}}}
\newcommand{\SpecialStringTok}[1]{\textcolor[rgb]{0.31,0.60,0.02}{{#1}}}
\newcommand{\ImportTok}[1]{{#1}}
\newcommand{\CommentTok}[1]{\textcolor[rgb]{0.56,0.35,0.01}{\textit{{#1}}}}
\newcommand{\DocumentationTok}[1]{\textcolor[rgb]{0.56,0.35,0.01}{\textbf{\textit{{#1}}}}}
\newcommand{\AnnotationTok}[1]{\textcolor[rgb]{0.56,0.35,0.01}{\textbf{\textit{{#1}}}}}
\newcommand{\CommentVarTok}[1]{\textcolor[rgb]{0.56,0.35,0.01}{\textbf{\textit{{#1}}}}}
\newcommand{\OtherTok}[1]{\textcolor[rgb]{0.56,0.35,0.01}{{#1}}}
\newcommand{\FunctionTok}[1]{\textcolor[rgb]{0.00,0.00,0.00}{{#1}}}
\newcommand{\VariableTok}[1]{\textcolor[rgb]{0.00,0.00,0.00}{{#1}}}
\newcommand{\ControlFlowTok}[1]{\textcolor[rgb]{0.13,0.29,0.53}{\textbf{{#1}}}}
\newcommand{\OperatorTok}[1]{\textcolor[rgb]{0.81,0.36,0.00}{\textbf{{#1}}}}
\newcommand{\BuiltInTok}[1]{{#1}}
\newcommand{\ExtensionTok}[1]{{#1}}
\newcommand{\PreprocessorTok}[1]{\textcolor[rgb]{0.56,0.35,0.01}{\textit{{#1}}}}
\newcommand{\AttributeTok}[1]{\textcolor[rgb]{0.77,0.63,0.00}{{#1}}}
\newcommand{\RegionMarkerTok}[1]{{#1}}
\newcommand{\InformationTok}[1]{\textcolor[rgb]{0.56,0.35,0.01}{\textbf{\textit{{#1}}}}}
\newcommand{\WarningTok}[1]{\textcolor[rgb]{0.56,0.35,0.01}{\textbf{\textit{{#1}}}}}
\newcommand{\AlertTok}[1]{\textcolor[rgb]{0.94,0.16,0.16}{{#1}}}
\newcommand{\ErrorTok}[1]{\textcolor[rgb]{0.64,0.00,0.00}{\textbf{{#1}}}}
\newcommand{\NormalTok}[1]{{#1}}

% cslreferences environment required by pandoc > 2.7

% Polyglossia
\usepackage{polyglossia}
\setmainlanguage{en-GB}

% localized quotes
\usepackage[strict,autostyle]{csquotes}

% Figures
\usepackage{graphicx,grffile}
\makeatletter
\def\maxwidth{\ifdim\Gin@nat@width>\linewidth\linewidth\else\Gin@nat@width\fi}
\def\maxheight{\ifdim\Gin@nat@height>\textheight0.8\textheight\else\Gin@nat@height\fi}
\makeatother
% Scale images if necessary, so that they will not overflow the page
% margins by default, and it is still possible to overwrite the defaults
% using explicit options in \includegraphics[width, height, ...]{}
\setkeys{Gin}{width=\maxwidth,height=\maxheight,keepaspectratio}

% Additional packages
\usepackage{natbib}             % Advanced Bibliography (citep...).
\usepackage{amsmath,amsfonts,amssymb}
\usepackage{breqn}              % Line breaks in equations
\usepackage{url}                % Line breaks in url's
\usepackage{enumitem}           % Line spacing in lists
  \setlist[itemize]{noitemsep,nolistsep}
  \setlist[enumerate]{noitemsep,nolistsep}
\usepackage{orcidlink}          % ORCID

% Tables
\usepackage{longtable,booktabs}
\usepackage{caption}
% These lines are needed to make table captions work with longtable:
\makeatletter
\def\fnum@table{\tablename~\thetable}
\makeatother
% longtable 2 columns
% https://tex.stackexchange.com/questions/161431/how-to-solve-longtable-is-not-in-1-column-mode-error
\makeatletter
\let\oldlt\longtable
\let\endoldlt\endlongtable
\def\longtable{\@ifnextchar[\longtable@i \longtable@ii}
\def\longtable@i[#1]{\begin{figure}[t]
\onecolumn
\begin{minipage}{0.5\textwidth}\scriptsize
\oldlt[#1]
}
\def\longtable@ii{\begin{figure}[t]
\onecolumn
\begin{minipage}{0.5\textwidth}\scriptsize
\oldlt
}
\def\endlongtable{\endoldlt
\end{minipage}
\twocolumn
\end{figure}}
\makeatother

% Full-width tables
\usepackage{tabu}
\renewenvironment{table}{\begin{table*}}{\end{table*}\ignorespacesafterend}

% Prevent overfull lines
\setlength{\emergencystretch}{3em}


% hyperref comes last
\usepackage{hyperref}           % Hypertext links, PDF bookmarks
  \hypersetup{%
    urlcolor=blue,%
    linkcolor=black,citecolor=black,colorlinks=true%
  }

%----------------------------------------------------------------------------------------
%	ARTICLE INFORMATION
%----------------------------------------------------------------------------------------

\JournalInfo{Version: hal-05512154v1} % Journal information
\Archive{juin 06, 2026} % Additional notes (e.g. copyright, DOI, review/research article)

\PaperTitle{Computation of Large Spatial Datasets with the~M~function} % Article title

\Authors{
Eric Marcon\textsuperscript{1 \orcidlink{0000-0002-5249-321X}} Florence Puech\textsuperscript{2 \orcidlink{0000-0002-5279-6878}}
} % Authors
\affiliation{
\textsuperscript{1}AgroParisTech, UMR AMAP, CIRAD, CNRS, INRAE, IRD, Univ Montpellier, Montpellier, France.\\ \hspace{1em} \\\textsuperscript{2}Université Paris-Saclay, INRAE, AgroParisTech, Paris-Saclay Applied Economics, F-91120 Palaiseau, France.\\ \hspace{1em} 
}
\affiliation{*\textbf{Corresponding author}: \href{mailto:eric.marcon@agroparistech.fr}{\nolinkurl{eric.marcon@agroparistech.fr}}, } % Corresponding author

\Keywords{Distance-based method, M-function, Performance test, R Package dbmss} % Keywords
\newcommand{\keywordname}{Keywords} % Defines the keywords heading name

%----------------------------------------------------------------------------------------
%	ABSTRACT
%----------------------------------------------------------------------------------------

\Abstract{
Increasing access to large geo-referenced datasets, coupled with the development of computing power, has encouraged the search for suitable spatial statistical tools. Distance-based methods have been extensively developed in several scientific fields to detect spatial concentration, dispersion or independence of entities at any distance and without any bias. Recently, \citet{Tidu2023} highlighted the qualities of Marcon and Puech's \emph{M} function, a relative distance-based measure, and also expressed reservations about the computation time required. Herein, we propose a methodology that specifies the processing of large spatialized datasets with the \emph{M} function using R software. The deterioration extent of the \emph{M} results was estimated and discussed as the gains it provides in computation time. We provided evidence that the individual location approximation generated information loss at substantially small distances, implying a trade-off between the smallest distance at which spatial interactions could be detected and computing performance. The R code used in the article is given for the reproducibility of our results. The computational performance of \emph{M} was conducted using two methods: (i) a precise evaluation of the computational time and memory requirements for geo-referenced data was conducted using the \emph{dbmss} package in R via performance tests, and (ii) based on \citet{Tidu2023}, we considered an approximation of the geographical positions of the entities.
}

%----------------------------------------------------------------------------------------

\begin{document}

\selectlanguage{en-GB}

\flushbottom % Makes all text pages the same height

\maketitle % Print the title and abstract box

\tableofcontents % Print the contents section

\thispagestyle{empty} % Removes page numbering from the first page

%----------------------------------------------------------------------------------------
%	ARTICLE CONTENTS
%----------------------------------------------------------------------------------------

\section*{Introduction}\label{introduction}
\addcontentsline{toc}{section}{Introduction}

Agglomeration of people and economic activities is a key topic in economics \citep{Fujita1999, Fujita2013, Glaeser2010}.
Because spatial concentration offers productivity gains due to cost reductions, economists and geographers in regional science have extensively studied clusters of economic activity, their consequences on territories and the importance of the determinants that explain them \citep{Marshall1890, Duranton2004}.
Policymakers are also particularly interested in the agglomeration of activities, since a better understanding of the locations and extent of the agglomeration of activities allows them to develop more effective economic policies addressing regional growth disparities and the measures needed to promote economic convergence or greater territorial cohesion \citep[see for example][]{Baldwin2003}.

Since industrial agglomeration is a core question in regional science, spatial concentration measures are widely employed to evaluate the agglomeration of activities \citep{Combes2004}.
Increasing access to large spatial datasets and computing power has encouraged the development of statistical analysis tools for processing such data most efficiently \citep{Baddeley2016}.
Empirical studies at highly detailed geographical levels have been proposed in recent years \citep[see for example][ or \citet{Baragwanath2026} based on satellite imagery]{Daspremont2025}.
In particular, the detection of spatial structures (attraction, repulsion, and independence) of individual spatialized data using analyses that are no longer based on zoned data but on geo-located data has received an increasing attention.
Such an approach has the advantage of preserving the exact positions of the entities analyzed contrarily to data aggregated at a given spatial level (for example regions, counties etc.).
Since the 2000s, a series of spatial concentration measures based on distances between geo-localized entities studied has been proposed (see \citet{Marcon2017} for a complete review).
By considering space as continuous, the advantage of these distance-based methods is to circumvent any statistical bias associated to a discretization of space into separate units, known as the MAUP-Modifiable Areal Unit Problem: see \citet{Openshaw1979} or \citet{Arbia1989}. Numerous studies have shown the importance of using such a methodology in social sciences \citep{Arbia1989, Marcon2003, Sweeney2022} or in exact sciences \citep{Cressie1993, Lentz2011, Dray2021}.
For example, \citet{Giuliani2014} published in their review an analysis based on continuous space on the technology manufacturing industry in the metropolitan areas of Milan and Turin in Italy.
They showed all the importance of taking into account the exact geographic position of the plants to provide a complete description of the location patterns of high- and medium-high-technology manufacturing.
For example, they detected that these firms were spatially aggregated at distances greater than 500\textasciitilde meters in the metropolitan area of Milan.
This result encourages developments that explain such agglomeration patterns.

One of the existing distance-based methods has been recently pointed out by \citet{Tidu2023}: the \emph{M}~function, a particular statistical measure proposed by \citet{Marcon2010}.
This distance-based method, hereafter referred to as \emph{M}, makes it possible to highlight spatial structures within a spatialized distribution (attraction, repulsion, and independence) from a study based on the distances separating the analyzed entities.
Various applications of \emph{M} have been proposed for detecting industrial agglomeration.
For example, \citet{Jensen2011} and \citet{Floch2018} analyze the spatial concentration different facilities in French cities.
At a regional level, \citet{Coll-Martinez2019} proposed an application of the \emph{M} function on creative industries on a Spanish metropolitan area while more recently, \citet{Canello2025} employed in particular \emph{M} to characterize the location patterns of footwear subcontractors in Italy.
The \emph{M} function has also been rapidly applied to other domains as such as the spatial distribution of ethnic groups for example \citep{Deurloo2008}.
However, while this measure preserves all the richness of individual geo-located data, it requires a longer calculation time than other distance-based measures, as it is a relative measure \citep[ reviewed the advantages and limitations of a dozen existing distance-based measures]{Marcon2017}.
This might be a limit of the use of the \emph{M} function if datasets are very large.
\citet{Tidu2023} proposed decreasing \emph{M}~calculation times by introducing a voluntary positioning error for the analyzed entities.
For example, in their study, industrial establishments in Sardinia (Italy) were not located at their exact addresses but at the centroid of their municipality.
This repositioning reduces calculation times, as the number of possible distances between establishments is limited to the distances that separate the centroids of the municipalities.
This motivation is similar to that of \citet{Scholl2013} who proposed approximating the distances between pairs of entities by grouping them into classes to compute the \(K_d\)~function \citep{Duranton2005}, another distance-based measure.
The method of \citet{Scholl2013}, implemented in the \emph{dbmss} package \citep{Marcon2014} for R \citep{R}, provides a considerable gain in memory use, another critical aspect of processing large datasets, with minor loss of accuracy.
However, the information loss owing to the approximation of the location of objects should imply a loss of accuracy in the estimation of their interactions at the same scale, which needs to be assessed.

\textbf{Herein, we propose testing the effectiveness of \citet{Tidu2023}'s method and help researchers choose the appropriate method to characterize the spatial structure of substantially large datasets}.
The article provides three main contributions to the literature.
First, we show the advantages of estimating \emph{M} on datasets with an order of magnitude of 100,000 points or less.
On a personal computer, the computation time becomes excessive beyond that.
Second, we study the effect of the geographical approximation of the locations of the analyzed entities in terms of the deterioration in information that this approach creates.
Third, we show that an alternate method, i.e.~random thinning of the point pattern \citep{Sweeney2005, Arbia2017}, is more efficient.

The layout of the article is mentioned here.
The next two sections present the \emph{M} function and the necessary data generated for the tests.
Large point sets (in the order of several tens of thousands of points) that are either completely random or geographically concentrated are considered.
The third section details the use of the \emph{dbmss} package to calculate~\emph{M} and its confidence interval from a table that gives the position and characteristics of the points or a matrix of distances between them.
The fourth section measures the performance of \emph{dbmss} as a function of the size of the set of points, in terms of computing time and memory requirements.
The fifth section tests the spatial approximation which involves grouping them at the center of the cells of a grid, following the approach of \citet{Tidu2023}, which positions them at the center of the administrative units of their location.
In the last section, we conclude and discuss the advantages and the limits of an approximation of the locations on the accuracy of the results and the computing time.

\section{Spatial concentration measures in regional science}\label{spatial-concentration-measures-in-regional-science}

\subsection{Spatial concentration measures and regional science}\label{spatial-concentration-measures-and-regional-science}

In regional science, spatial concentration measures are widely employed to evaluate the agglomeration of activities \citep{Combes2004}.
These statistical measures provide a simple way to describe the geographic distribution of industrial activities across regions.
Industrial agglomeration generates productivity gains for firms and benefits for households \citep{Fujita1999, Glaeser2010, Fujita2013} and thus provokes growth disparities between territories.
Regional economic inequalities can be easily detected and quantified with spatial concentration measures.
In that way, a lot of researchers have employed spatial concentration measures for a better description and understanding of industrial clusters \citep[see][ among others]{Feser2000, Giuliani2014, Kerr2015}.
Moreover, since spatial concentration measures are proven to be a relevant technique to describe agglomeration across territories, some studies rely on these spatial concentration estimates to explain the mechanisms that generate the industrial clustering.
The pioneer explanation of the agglomeration forces is due to \citet{Marshall1890} who identifies three main determinants of the industrial spatial concentration: a specialized labor pooling, the proximity of producers of intermediate goods (input sharing), and the opportunity to benefit from technological externalities (knowledge spillovers).
These factors explain why firms concentrate their activities in a limited number of specific geographic areas: to reduce their costs and risks associated to their production.
These three determinants are still widely accepted today, even though the modern view proposed by \citet{Duranton2004} of these economic mechanisms is approached by a better matching, a better sharing, and a better learning for firms and households.
A series of articles have empirically evaluated the importance of the determinants of agglomeration by regressing estimates of spatial concentration measures on the main factors that generate agglomeration.
This technique constitutes a very simple way to quantify the importance of all of the determinants of industrial agglomeration: see for example the empirical studies of \citet{Rosenthal2001} or \citet{Ellison1999}, or the review of \citet{Rosenthal2004}. Estimating at best the importance of all factors at work is straightforward, notably for policymakers that provide effective economic policies in order to foster regional growth and to attenuate disparities between regions.

\subsection{Recent developments in spatial concentration measures}\label{recent-developments-in-spatial-concentration-measures}

Based on the importance of the industrial geographic concentration, a great deal of attention has been paid in economics and in geography to develop statistical measures that are able to describe the spatial concentration precisely and without any statistical bias \citep{Arbia2021}.
Today, many indices coexist in the literature but two distinct groups are {[}F1.1{]} clearly identified depending both on the nature of the data used and on their statistical treatment.

The first group of measures divides space into a given number of exclusive zones, for example data are aggregated at the county level or the regional one.
The well-known and widely used location quotient \citep{Florence1972}, the \citet{Gini1912} index or the \citet{Ellison1997} index belong to that group.
It is out of the scope of this article to present in details these indices but a discussion of these spatial concentration measures is given in \citet{Sweeney2004}, \citet{Bickenbach2008} or \citet{Combes2008}.
Entropy measures also belong to that first group, even if they are less employed in spatial economics or in quantitative geography \citep[see for example][ for an evaluation of spatial concentration of activities across regions with the \citet{Theil1967} index]{Brulhart2005}. The common advantage of all of these measures is the simplest way to compute them: they do not need intensive calculations, data are spatially aggregated so they are easy to find and compute. For example, \citet{Marcon2019} employed European data (available at \url{https://ec.europa.eu/eurostat/web/regions-and-cities}) to measure industrial employment disparities across European regions. The shared limit of all of the previous indices presented is that they rest on a pre-defined zoning of space. As a consequence, all the information lower than the zoning delimitation is simply lost. This is far from being ideal for example for researchers that want to study the very localized interactions between firms, the face-to-face contacts between individuals etc{[}F2.1{]} .
In addition, a series of studies put in light that all indices based on a discrete space are exposed to some statistical bias known as the Modifiable Areal Unit Problem-MAUP \citep{Openshaw1979, Arbia1989}.
For example, it has been proven that the position of the frontiers of the zones or the geographical scale chosen for the analysis alter the findings on the spatial concentration levels \citep{Arbia2001}.
This is important at both for the precision of the agglomeration levels but also to find the \enquote{right} scale to analyze agglomeration factors \citep{Rosenthal2020}.

An answer the MAUP issues was proposed by the development of a second group of measures that not divide space into a given number of zones.
This second group encompasses all \enquote{distance-based measures} (DBMs) based on point pattern analysis \citep{Marcon2017}.
In that case, the exact geographic position of firms is preserved and the spatial concentration levels of agglomeration is calculated on the distances between the geo-located firms studied.
By preserving the exact positions of the firms analyzed, DBMs circumvent all statistical bias associated to a discretization of space into separate units.
By considering space as continuous, DBMs evaluate the spatial concentration for all possible distances.
One of the great advantages of all distance-based methods is to detect the exact distance at which spatial concentration (in case of attraction) or dispersion (in case of repulsion) occur.
In the last decades, many developments have been proposed in geography and in economics to be able to grasp any definition of the spatial concentration analyzed (different benchmarks distributions, weighting or not the plants by their number of employees etc.).
A series of spatial concentration measures based on distances have been proposed \citep[ proposed a typology of around a dozen of existing DBMs]{Marcon2017}. Among the first applications of distance-based methods, the works of \citet{Barff1987} or \citet{Sweeney1998} are of particular interest. They investigated the relationship between the agglomeration of manufacturing activities and the characteristics of the plants (capital intensity, size, etc.). \citet{Sweeney1998} notably gave support to an existing inverted U-shaped relation between the spatial manufacturing concentration and the size of plants in North Carolina (USA).{[}F3.1{]} Today, DBMs face a challenge: they must be able to address the huge amount of micro-geographic data that are increasingly available for researchers \citep{Piacentino2021}.
For example, dealing with satellite data \citep[\citet{Daspremont2025}, \citet{Baragwanath2026}]{Bachtrogleunger2023} is an opportunity to provide detailed locational analysis of the economic activities.
This kind of data opens new research questions both on the computational capacity to manage very large datasets and the possible reductions of computing time.
In a recent paper \citet{Tidu2023} focus their attention on one of the existing distance-based methods, the \emph{M} function \citep{Marcon2010}, that presents many advantages for analyzing the location of industries.
In the next sections, we shall discuss the computing performances of the \emph{M} function and we also try to appreciate its capacity to deal with very large datasets.

\section{\texorpdfstring{The \emph{M} function}{The M function}}\label{the-m-function}

\subsection{Main idea}\label{main-idea}

\citet{Marcon2010} introduced the \emph{M}~function to evaluate the dependence between geo-located points without relying on a specific zoning of space.
Similar to any distance-based method, the calculation of~\emph{M} is based on distances that separate entities studied (establishments, shops\ldots).
The idea of~\emph{M} is simple: it compares two proportions of neighbors of interest, a local one to a global one.
The local one is defined as the proportion of neighbors of interest within a distance~\emph{r}.
The global one is the same proportion but defined on the entire territory.
This comparison of ratios enables the detection of:

\begin{itemize}
\tightlist
\item
  spatial concentration (attraction) of entities if the proportion of local neighbors is greater than the one observed on the entire territory,
\item
  spatial dispersion (repulsion) of entities if the relative proportion of local neighbors is lower than the one observed all over the territory,
\item
  independence between entities if the local distribution of neighbors does not differ from the global one.
\end{itemize}

This comparison of proportions of neighbors defines~\emph{M} strictly as a \emph{relative} distance-based measure \citep{Marcon2017}.
The term \emph{topographic} distance-based measures is preferred for those that use the surface area as a benchmark, as the well-known Ripley's \emph{K} function or \emph{L} function \citep{Ripley1976, Ripley1977, Besag1977}.
The \emph{M}~function is also defined as a \emph{cumulative} distance-based method because the local environment is appraised within a distance \emph{r} rather than at a distance~\emph{r}.
The possibility to detect exactly at which distance(s) the spatial concentration or dispersion appears coupled with the interpretation of results opens the way to precisely describe the distribution of entities under study.
An easy-computation of~\emph{M} is possible owing to the \emph{dbmss} R~package \citep{Marcon2014}.

\emph{M}~was first introduced in the field of economics and geography.
\citet{Marcon2010} proved that this function satisfies all the requirements of \citet{Duranton2005} for the evaluation of the spatial distribution of industries.
Since its introduction, various studies have described the spatial locations of industries by using~\emph{M}; for example, \citet{Jensen2011} studied the location of shops at an urban level, \citet{Coll-Martinez2019} analyzed creative industries at a metropolitan level.
This methodology has also been rapidly applied in other domains including biology \citep{Fernandez-Gonzalez2005}, geography \citep{Deurloo2008}, ecology \citep{Marcon2012}, and seismology \citep{Nissi2013}.
The \emph{M}~function is now included in general textbooks of spatial statistics such as \citet{Arbia2021} but it is less popular than the \(K_d\) function of \citet{Duranton2005}; see \citet{Chain2019}.

\subsection{Definition}\label{definition}

The \(M\) function is based on the point process theory, in line with the distance-based method proposed by Ripley and Besag distance-based method \citep{Ripley1976, Ripley1977, Besag1977}.
The \(M\) function was developed to analyze interactions among entities in a heterogeneous space.
In other words, within this statistical framework, we consider that any entity analyzed does not have the same probability to locate everywhere on the territory (the \emph{first-order property of the point pattern} is its intensity).
Subsequently, upon controlling for space heterogeneity, we can identify interactions and detect spatial concentration or dispersion (the \emph{second-order property of the point pattern}).
Space heterogeneity is a consistent assumption for studying the agglomeration of industries \citep[refer to the discussion of][ on that subject]{Duranton2005}.

The definition of the intra-type version of~\emph{M} is as follows.
It compares the relative proportion of entities of interest up to each distance~\(r\) to the same ratio but defined over the entire territory under study.
In this article, we only consider the intra-type version of~\emph{M}: we study the spatial structure of neighboring points of the same type (called \emph{points of interest}) as the points at the centers of the disks of radius~\(r\).
In mathematical terms, we denote the terms mentioned below:

\begin{itemize}
\tightlist
\item
  \(x^s_i\), the location of point~\(i\) of the reference type~\(s\), at the center of the disk (the point whose neighborhood is to be analyzed).
\item
  \(x^s_j\), the location of a neighbor \(j\) of the same type as point~\(i\).
\item
  \(x_j\), the location of a neighbor~\(j\) of~\(i\), regardless of its type.
\item
  \(w(.)\), the weight of a given neighbor. In that sense, \(w(x_j)\) defines the weight of a neighbor~\(j\) of~\(i\).
\item
  \(W_s\), the total weight of the points~\(x^s_j\).
\item
  \(W\), the total weight of all points of the dataset, regardless of their type.
\item
  \({\mathbf 1} \left( \left\| x^s_i - x_j \right\| \le r \right)\), the indicator function is equal to one if \(x_j\)~is in the neighborhood of \(x^s_i\), e.g., the distance between \(x^s_i\) and \(x_j\) is at most equal to~\(r\), zero otherwise.
\end{itemize}

The intra-type \(M\) function is defined as \eqref{eq:M}.

\begin{equation} 
  \hat{M}\left(r\right) = {\sum_i{\frac{\sum_{j \ne i}{{\mathbf 1} \left( \left\| x^s_i - x^s_j \right\| \le r \right) w \left( x^s_j \right)}}{\sum_{j \ne i}{{\mathbf 1} \left( \left\| x^s_i - x_j \right\| \le r \right) w \left( x_j \right)}}}} / {\sum_i{\frac{W_s - w \left( x^s_i \right)}{W - w \left( x^s_i \right)}}}
\label{eq:M}
\end{equation}

Several remarks must be made.
The first remark is that the benchmark value of~\emph{M} is equal to one, regardless of the distance considered.
It means that for any radius~\emph{r}:

\begin{itemize}
\tightlist
\item
  If the estimated \(M\) result is above one, the local value of the ratio is greater than the global one: a spatial concentration of entities of type \(s\) within that radius is thus detected.
\item
  If the estimated \(M\) result is under one, the local value of the ratio is lower than the global one: a spatial dispersion of entities of type \(s\) within that radius is thus detected.
\end{itemize}

The second remark concerns the significance of the results.
A confidence interval can be generated using Monte Carlo simulations following \citet{Marcon2010}.
A risk level is chosen (for example 5\%) as well as the number of simulations.
The greater the number of simulations, the longer is the duration of the calculation of \emph{M}.
Third, the package \emph{dbmss} \citep{Marcon2014} on the R software \citep{R} can be used to compute the \(M\) function.
In most studies, the Euclidean distance is preferred to calculate~\emph{M}, but the \emph{dbmss} package can also be used for network distances.
We discuss that point hereinafter for large datasets.

\subsection{Simulated data}\label{simulated-data}

The datasets we consider in this article were obtained by simulation.
The R code is given in the appendix, which allows perfect reproducibility of the examples treated.

\subsubsection{Drawing the points}\label{drawing-the-points}

A set of points is drawn using the Poisson process (whose expectation of the number of points is 100,000) in a square window of side one.
Each point is assigned a qualitative mark: \enquote{Case} or \enquote{Control}.
95\% of points are Controls.
Additionally, 5\% are Cases, whose spatial structure is studied.
The weight of the points is drawn from a gamma distribution with free shape and scale parameters.

In this example, the drawing of Controls points is completely random (\emph{complete spatial randomness}: CSR), i.e., there is no simulation of attraction or dispersion.
On the contrary, the spatial distribution of Cases is aggregated, which means that Cases are spatially concentrated (like clusters).
Practically, sets of aggregated points are drawn in a \citet{Matern1960} process.



\scriptsize

\begin{figure}

{\centering \includegraphics[width=1\linewidth]{Rev1_files/figure-latex/XMaternFig-1} 

}

\caption{Random drawing of a set of points where the Cases (red) are aggregated and the Controls (blue) are distributed completely randomly. The size of the points is proportional to their weight.}\label{fig:XMaternFig}
\end{figure}

\normalsize

Cases shown in Figure~\ref{fig:XMaternFig} indicate visible aggregates.
Controls are distributed completely randomly.
A careful analysis of Figure~\ref{fig:XMaternFig} shows a limited number of tiny white spaces on the square window, indicating locations with no points (whatever the type).
The scarcity of empty spaces is because of the high number of simulated points.

\subsubsection{Gridding the space}\label{gridding-the-space}

The simulation of the Cases obtained by the Matérn process is considered and the window is split into a 20~x~20 square grid.
This partition simulates the approximation of the position of the points of an administrative unit to the position of its center.
Moreover, this grid size is consistent with that of \citet{Tidu2023}, which facilitates a comparison of our results.
The choice of the optimal level of the grid remains an open question, as \citet{Arbia2021} noticed (p.109): \enquote{\emph{Unfortunately, the choice of the partitioning scheme is usually arbitrary and an optimal criterion to guide this choice is not available}.}

The approximated position of points is depicted on the map presented in Figure~\ref{fig:GroupedFig}.
Each cell comprises only one point of each type, whose weight is the sum of the weights of the individual points.



\scriptsize

\begin{figure}

{\centering \includegraphics[width=1\linewidth]{Rev1_files/figure-latex/GroupedFig-1} 

}

\caption{Repositioning of points in an arbitrary grid. The absence of Cases in a cell is easily detected (single-color blue dot), as is the strong presence of Cases in a cell (two-color dot, but predominantly red).}\label{fig:GroupedFig}
\end{figure}

\normalsize

\emph{M} values can be calculated from the original point set or its approximation.

\subsection{\texorpdfstring{Computing \emph{M} using the \emph{dbmss} package}{Computing M using the dbmss package}}\label{computing-m-using-the-dbmss-package}

In the \emph{dbmss} package, the data are a set of points or a distance matrix.
The set of points in Figure~\ref{fig:XMaternFig} is used.

\subsubsection{Point pattern}\label{point-pattern}

The \texttt{Mhat()} function in the \emph{dbmss} package is used to estimate the \emph{M} function.
The theoretical reference value for \emph{M} is one, as this function relates the proportion of Cases up to a distance \(r\) to that observed across the entire window.
The aggregation of Cases will be highlighted by values of \emph{M}~\textgreater~1 (the relative presence of Cases is greater locally than across the entire window) and the dispersion of Cases by values \textless~1.
The Euclidean distance is used to estimate the distance between points.

Figure~\ref{fig:MFig} shows that \emph{M} detects an agglomeration of Cases, which is in line with the simulation of this type of point (Controls having a completely random location on the window).
The advantage of a function based on distances is clearly visible: for any distance, the level of spatial concentration is precisely estimated.
It enables the detection of distances at which the attraction phenomena occur and are the most important (for functions whose values can be compared at different radii, such as \emph{M}).
In addition to estimating the \emph{M} function, the \texttt{Menvelope()} function can be used to calculate its global confidence interval \citep{Duranton2005} under the null hypothesis of random point location.
The confidence interval of the null hypothesis is centered on one, and is narrow because of the large dataset.
Lastly, the necessary simulations can be parallelized to save time.



\scriptsize

\begin{figure}

{\centering \includegraphics[width=1\linewidth]{Rev1_files/figure-latex/MFig-1} 

}

\caption{Value of \emph{M} as a function of the distance from the reference point. The 95\% confidence envelope, obtained from 100 simulations, appears in gray and is centered on the value of one.}\label{fig:MFig}
\end{figure}

\normalsize

\subsubsection{Distance matrix}\label{distance-matrix}

Matrices can be used to process non-Euclidean distances (transport time, road distance, etc.) which cannot be represented by a set of points.
The \texttt{Mhat()} and \texttt{MEnvelope()} functions are the same, and provide the same results irrespective of the the data form (point set or distance matrix).
Details are given in the appendix.

The size of distance matrices is the square of their number of points, i.e.~10\textsuperscript{10} cells for 100,000 points.
R does not handle them because it relies on integer values to index them: square matrices cannot be larger than 46340 rows and columns, i.e., the square root of the largest supported integer value.
This does not allow dealing with large datasets.
Thus, this approach is not explored further in this paper.

\subsection{Computational performance}\label{computational-performance}

The use of \emph{M} to characterize the spatial structure of large sets of points might be limited by its computing time or the memory required.
In this section, we investigate these two potential limitations.

\subsubsection{Computing time}\label{computing-time}

The distances between all pairs of points must be calculated to estimate~\emph{M}.
Therefore, it is expected that the calculation time will increase as the square of the number of points.
The time required for the exact calculation is evaluated for a range of numbers of points (Figure~\ref{fig:TestTimeMemFig}a).

\textbf{The calculation time is related to the size of the set of points by a power law}.
It increases less rapidly than the square of the number of points.
It can be estimated precisely (\(R^2=\) 0.98) using the mentioned relation: \(t=t_0 (n/n_o)^p\), where \(t\) denotes the average time for \(n\)~points (e.g., 4.69 seconds for 100,000 points), knowing the time \(t_0\) for \(n_0\) points, and \(p\) is the power relation (here: 1.6).



\scriptsize

\begin{figure*}

{\centering \includegraphics[width=1\linewidth,]{Rev1_files/figure-latex/TestTimeMemFig-1} 

}

\caption{(a) Calculation time (seconds) and (b) memory required (MB) to estimate~\emph{M} as a function of the size of the set of points. The measures were repeated 10 times. The bars represent the~+/-~1 standard deviation interval.}\label{fig:TestTimeMemFig}
\end{figure*}

\normalsize

\subsubsection{Memory}\label{memory}

The memory used is evaluated for the same data sizes (Figure~\ref{fig:TestTimeMemFig}b).
\textbf{The memory required increases linearly with the number of points and is never critical for point set sizes that can be processed within reasonable times}.
This highlights \citet{Tidu2023}'s conclusion regarding the power and computation time required when using~\emph{M} on large datasets.

\section{A guide to dealing with large datasets}\label{a-guide-to-dealing-with-large-datasets}

Computation time becomes critical when very large datasets are considered, i.e.~above 100,000 points as shown in the previous section.
The M function is often used to test a point pattern against a null hypothesis, requiring simulations of a null point pattern and the same computing time for each of them, hence multiplying the computing time.

We explore here two possible ways to reduce the computational burden: approximating the position of points, and thinning the point pattern.

\subsection{Effects of approximating the position of points}\label{effects-of-approximating-the-position-of-points}

\subsubsection{Two main effects}\label{two-main-effects}

Unambiguously, approximating the position of the points generates a loss of information: in each grid cell, the distance between all the points is set to zero, and the distance between two points in different cells is approximated using the distance between the centroids of the two cells.
The first consequence is that \textbf{no spatial structure in each cell can be detected because the information is completely lost.} The merits of any distance-based methods are detecting the exact geographic scale(s) where any patterns of aggregation, dispersion or independence of points exist.
Grouping points at the centroid of cells erases any spatial structure under the size grid.
The same limit under the size grid is faced for distributions with multiple patterns of points (that is a coexistence of attraction and dispersion of points depending on the distance considered).
\citet{Tidu2023} gathered an average of 320 points in every 377 municipalities of Sardinia.
However, the local structure of these points was overlooked.
Moreover, we suspect \textbf{bias in~\emph{M} estimation based on a scale of the order of the magnitude of the size of the cells}, which should decrease with distance, when the relative size of the cells becomes negligible.

\subsubsection{Test setting}\label{test-setting}

The effect of the location approximation is tested on a set of aggregated points, similar to the real \citet{Tidu2023} data, and on a completely random point pattern.
Both comprise 100,000 points such that groups include 250 points (100,000 points divided by 400 groups) on average.

20 sets of completely random and of aggregated points (100,000 points with 5\% of Cases) were simulated.
The exact calculation and the calculation on the grid points were performed on each set of points to evaluate the effect of the approximation.



\scriptsize

\begin{figure*}

{\centering \includegraphics[width=1\linewidth,]{Rev1_files/figure-latex/MapproxFig-1} 

}

\caption{Average estimate of~\emph{M} from the exact position of the points compared with the values obtained by grouping the points for CSR (a) and Matérn (b) point patterns.}\label{fig:MapproxFig}
\end{figure*}

\normalsize

\subsubsection{Results}\label{results}

The mean values of the estimates of~\emph{M} are presented in Figure~\ref{fig:MapproxFig}.
The size of the grid cells is equal to 0.05.
All neighbors at distances less than this threshold are placed at zero distance.
The estimate of the corresponding~\emph{M} function (Grouped~\emph{M}) is constant in that case, up to this threshold.
When the actual point pattern is not structured, some artifactual aggregation is generated at a small scale.
In contrast, the local aggregation of the Matérn pattern is underestimated.
Figure~\ref{fig:MapproxFig} shows that at substantially short distances the grouped~\emph{M} plot is below the exact~\emph{M} plot obtained using the actual position of points.

\citet{Tidu2023} tested the effect of grouping the points by the correlation between exact and approximated \emph{M} values.
We refrained from following them because the main source of variation of the grouped-point \emph{M} values is that of the grouping itself: when the number of points per group is small, groups considerably vary between simulations, and the correlation is weak.
It increases with the size of the data (an illustration is given in the appendix).
A substantially high correlation does not exclude systematic errors.
Therefore it is not considered an appropriate statistic here.



\scriptsize

\begin{figure}

{\centering \includegraphics[width=1\linewidth]{Rev1_files/figure-latex/MCompareFig-1} 

}

\caption{p-values to reject the null hypothesis (H0) of independence of the point locations tested by the \emph{M} function. The colors of the curves represent CSR or aggregated point patterns. Solid lines provide the exact p-values of \emph{M}, and dotted lines provide the values estimated on grouped point patterns. The horizontal, dotted line corresponds to the significance threshold, i.e., 97.5\%, to reject H0 in a 5\% risk-level two-tailed test.}\label{fig:MCompareFig}
\end{figure}

\normalsize

We chose the mentioned test.
The \emph{M} function is basically used as a test against the independence of point locations.
To assess the impact of grouping the points on the test result, Figure~\ref{fig:MCompareFig} shows its average p-value, computed among 20 simulations of each point pattern.
At each distance, the~\emph{M} value is compared with that obtained with point patterns simulated based on the null hypothesis, which is rejected if the actual value is outside of the 95\% central quantiles of the simulations.

In the case of CSR, the p-value to reject independence around 50\%, regardless of the distance, and whether the points are grouped or not.
In the case of aggregated patterns, the null hypothesis is rejected correctly when the points are grouped.
At small scale, below the size of the cells, \emph{M} values are actually that of the cell size: the test is correct because the point process we used is aggregated at all scales.
The point pattern might have been repulsive at the small scale, but this information is destroyed by grouping the points: \emph{M} and their p-values are not reliable below the size of the cells.

\textbf{Since all information below the grid size is lost, the approximation might or might not be acceptable depending on the research question and the spatial scale at which interactions occur. Above the grid size, the \emph{M} function is less affected by the approximation and the gap becomes rapidly negligible}.

We can easily understand that applying a grid is an attractive technique to save calculation time, but the risk is missing the geographic scale at which the spatial phenomenon occurs.
Let us come back to a core question in regional sciences given in section 1, the industrial clustering, to illustrate the challenge of finding the optimal grid size.
Knowing the real extent of agglomeration economies is not an easy task but it has been proven that they decay rapidly as the distance increases \citep{Rosenthal2003, Rosenthal2020}.
Empirical studies do not always conclude on the same geographic scale on which plants agglomerate.
Many factors influence the findings: the industry considered, the distance-based measure employed, the area considered, etc.
To give some examples, \citet{Giuliani2014} found that plants in high- and medium-high-technology manufacturing sectors are spatially aggregated at distances greater than 500\textasciitilde meters in the metropolitan area of Milan while no significant pattern is observed in the same industry in the metropolitan area of Turin.
In a pioneer study, \citet{Sweeney2004} detected for Los Angeles that the highest level of spatial concentration of establishments in electronics appears within a distance of 1km while for textiles plants the agglomeration is the greatest within a distance of 4km.
In their seminal article, \citet{Duranton2005} showed that location patterns may be very different from an industry to another by taking four \emph{ad hoc} industries in the United Kingdom (basic pharmaceuticals, pharmaceutical preparations, other agricultural and forestry machinery, and machinery for textile, apparel and leather production).
These results are perfectly in line with \citet{Feser2002}'s findings for the United States.
The latest proposed a very detailed analysis for a large number of industries in fourteen American metropolitan areas (see tables on pages 245-246 of their article).
The discrepancies in the location patterns at both within and between industries are appreciable.
In sum, spatial proximity is a multiscale problem.
Knowing \emph{a priori} the optimal grid size is not easy but one thing is certain: missing the geographical scale at which agglomeration economies is very embarrassing while DBM are available and perfectly preserve the richness of the spatial information of geolocated datasets.

\subsection{Random thinning}\label{random-thinning}



\scriptsize

\begin{figure*}

{\centering \includegraphics[width=1\linewidth,]{Rev1_files/figure-latex/MthinFig-1} 

}

\caption{(a) Confidence envelope of the estimation of the M function after thinning a 1-million-point pattern with aggregated Cases to 10\% of its size. The mean M values among thinned patterns (dotted, red line) can't be distinguished from the exact values (black line). The grey envelope corresponds to the 95\% confidence level. (b) relative Root Mean Squared Error of the estimation due to thinning, i. e. the expected relative error of estimation, always below 2\%.}\label{fig:MthinFig}
\end{figure*}

\normalsize

Random thinning consists of reducing the number of points by applying the same probability of being kept to each point.
The resulting dataset is thus a representative sample of the actual one.

\citet{Sweeney2005} showed that the K and D functions, that are other DBMs, are invariant under random thinning.
So is the M function: it is obtained by dividing a number of neighbors of interest by a total number of neighbors.
Both terms are reduced in the same proportion by random thinning, keeping the ratio unchanged.

Compared to approximating the position of points, random thinning does not degrade the information of each point but reduces the amount of data.
The acceptability of the method can be tested rigorously by simulating randomly thinned point patterns and comparing the distribution of their M function to that of the original pattern.
A million-point pattern is drawn with aggregated Cases and thinned with down to 10\% of its size.
The M values of the original and 100 thinned patterns are compared.
Figure~\ref{fig:MthinFig}a shows the mean value and the confidence envelope of the M function applied to thinned patterns.
As expected, thinning causes no bias.
The relative Root Mean Squared Error (RMSE), that is the expectation of the error caused by thinning, is shown on Figure~\ref{fig:MthinFig}b.
Since the bias is zero, it is simply the coefficient of variation of the simulated M values at each distance.
It is very small, around 2\%, at short distances and negligible at large distances.

As \citet{Sweeney2005} already noted for the K function, the variance of the estimation of M increases when the number of points decreases.
Thinning a 100,000-point dataset to 1000 points results in a very fast and unbiased estimation of the M function but RMSE is well over 100\% at short distance (Figure in supplementary material).
Thinning is an efficient method to reduce huge, untractable datasets to large, computable ones, but is not an appropriate way to drastically reduce computing time by thinning large datasets to small ones.

\section{Discussion and conclusions}\label{discussion-and-conclusions}

\textbf{The computation burden of estimating~\emph{M} on large datasets might be a problem}.
The calculation time for \emph{M} is \textless~5 seconds for a set of 100,000 points on a modern computer\footnote{The results presented here were obtained on a GitHub-hosted Mac OS runner with a virtual 24-core 12th Gen Intel(R) Core(TM) i7-12850HX, similar to a fast laptop computer.} and requires 25 MB of RAM.
Therefore, calculating a confidence interval from 1,000 simulations requires less than 83 minutes.
For a set of five million points, the power law predicts around 36 minutes of computing time.
Accordingly, 1,000 simulations would take around 25 days.

Owing to parallelization, a calculation server would drastically increase performance, but at the cost of the complexity of implementation that limits its use.
If we limit ourselves to the computing power of a personal computer, \textbf{exact calculation is fully justified for data of the order of 100,000 points}: less than an hour is sufficient to calculate confidence intervals.
Since parallelizing the simulations is offered with no effort by the \emph{dbmss} package, this time can be reduced by a considerable factor based on the available hardware, e.g., by a factor of 2 to 6 using modern multicore central processing units.
Beyond that, approximating the location reduces the size of the set of points to the number of locations selected.
It may be up to 100,000 locations to keep the computing time acceptable, regardless of the size of the original dataset.
The size of the grid \citep[or the administrative scale of the aggregation of points in][]{Tidu2023} should be chosen based on the scale of the interactions under study: they cannot be characterized correctly at distances below it.

Regarding the errors generated in the estimates of~\emph{M} when the approximation of location is used, our findings support those of \citet{Tidu2023}, which mentions strong correlations between \emph{M} values computed from exact and approximated Italian company location data.
The spatial approximation problem originates from the loss of information regarding possible interactions at distances smaller than the size of the grid.
This situation needs to be analyzed attentively.

Moreover, \citet{Arbia2017} investigated applying distance-based methods to spatial datasets that include positional errors.
They proposed a first evaluation of the consequences of an \emph{\enquote{unintentional positional error}} owing to the uncertainty of the location of a part of the studied entities, e.g., when their address is not accurate.
In that scenario, these uncertain geo-localized entities are placed at the centroid of the zone considered, exactly as in \citet{Tidu2023}.
Arbia et al.~showed on a real case (Italian manufacturing firms) that the error measurement was less severe than expected.
Their explanation rests on the definition \emph{M}, a relative measure: a compensation effect of positional errors is suspected between the local and the global ratios.

\textbf{An approximation of the spatial locations might be considered to save computation time, given the strong agreement observed between \emph{M} values on exact and approximated data above the grid size, but the information related to interactions at smaller scales is lost. As we previously noticed, it is challenging to choose the optimal size of the grid, requiring a certain degree of caution in the use of the approximated locations}.
For example, careful use of location approximation is recommended if short-distance interactions are suspected, such as information externalities or contagion phenomena.
Choosing the approximation scale is a trade-off between accuracy, i.e., a small distance threshold above which results are accurate, and speed with a coarse grid.

\section*{Appendix}\label{appendix}
\addcontentsline{toc}{section}{Appendix}

R code is available at the following address: \url{https://ericmarcon.github.io/MLargeDataSets/Appendix.pdf}

\section*{Acknowledgments}\label{acknowledgments}
\addcontentsline{toc}{section}{Acknowledgments}

We thank the editor and three anonymous referees for helped us improve the article.
We also thank participants of the SEW 2025 and especially Vincenzo Nardelli for helpful advice.
Eric Marcon benefited from an \enquote*{Investissement d'Avenir} grant managed by the Agence Nationale de la Recherche (LABEX CEBA, ref. ANR-10-LBX-25) and Florence Puech gratefully acknowledges financial support from INRAE.

%----------------------------------------------------------------------------------------
%	REFERENCE LIST
%----------------------------------------------------------------------------------------

\bibliographystyle{chicago}
\makeatletter
% The filename has .bib extension that must be eliminated
\filename@parse{references.bib}
% parse stores the file name in base. Extension starts at the first dot, so don't use dots in file names.
\bibliography{\filename@base}
\makeatother


%----------------------------------------------------------------------------------------

\end{document}
